\documentclass[10pt,a4paper]{amsart}
\usepackage[margin=19mm]{geometry}
\usepackage{amsmath,amssymb,mathtools}
\usepackage[scr=rsfs,cal=euler]{mathalfa}
\usepackage{xfrac}
\usepackage[unicode,hidelinks,bookmarks=false]{hyperref}
\newcommand{\ie}{i.e.\ }
\newcommand{\cf}{cf.\ }
\newcommand{\resp}{respectively\ }
\newcommand{\Subsection}{\subsection}
\providecommand{\nobreakdash}{\nobreak\hbox{-}}
\setlength{\emergencystretch}{2em}
\multlinegap0pt
\usepackage{bm}
\usepackage{bbm}
\usepackage{dsfont}
\usepackage{xspace}
\usepackage{cancel}
\usepackage{mathtools}
\usepackage{amsthm}
\makeatletter
\@ifundefined{theorem}{\newtheorem{theorem}{Theorem}}{}
\makeatother
\newcommand{\Aut}{\mathop{\mathrm{Aut}}}
\title[On minimal graphs, fixing sets and base size sets for hamilt]{On minimal graphs, fixing sets and base size sets for hamiltonian groups}
\author{Kirti Sahu, Ranjit Mehatari}
\date{Source: arXiv:2601.21575v2 (CC BY 4.0)}
\begin{document}
\maketitle

\noindent\emph{Source and licence.} On minimal graphs, fixing sets and base size sets for hamiltonian groups. Version arXiv:2601.21575v2 (CC BY 4.0), distributed under the Creative Commons Attribution 4.0 International licence, as recorded in the arXiv licence metadata for this exact version and shown on the abstract page https://arxiv.org/abs/2601.21575v2. Full rights notice: licenses/arxiv-2601.21575-CC-BY-4.0.txt.\par\smallskip

\noindent\textit{Editorial scope.} Counted unit: the Theorem whose source label is \texttt{ham\_th1} in the article (source0.tex lines 183--186, source label \texttt{ham\_th1}), delivered together with the article's own complete original proof of it (source0.tex lines 188--268, one \texttt{proof} environment of 81 lines). No mathematics is written, repaired or re-derived: statement and proof are the article's own text line for line. Required context retained uncounted: none. Every definition, display and label of those lines is retained unchanged. Omitted from this package and not redistributed: 1--182, 187--187, 269--573. Nothing inside a retained excerpt is omitted or altered: every retained body is delivered exactly as the article's lines are written, so no omission receipt of any kind accompanies this package. Citations: the retained excerpts carry no citation command at all, so no bibliography entry of the article is transcribed and no reference is redistributed. Reference adaptation: none was needed and none was made; every reference inside the delivered unit resolves inside it, so no unresolved reference or citation is produced. Printed slips kept literal: no printed word, symbol, hypothesis or number of the counted statement or of its proof was altered. Layout and class: the article's own class is \texttt{article} with packages including amsmath, amsthm, fullpage, amssymb, epsfig, tikz, graphicx, url, authblk; the wrapper is the lane's pinned \texttt{amsart} editorial wrapper at 10pt on A4, which restates the theorem-like environments the retained text needs and adds the article's own macro definitions that the retained text needs (\texttt{\textbackslash Aut}), with extra wrapper packages (bm, bbm, dsfont, xspace, cancel, mathtools). The delivered numbering is the unit's own. Assets: no retained span contains a figure environment, a graphics-inclusion command, a PDF-page inclusion, a table or a TikZ picture; no image file is copied into this package, the package declares no asset dependency and no image file is redistributed with it. Licence: version arXiv:2601.21575v2 of this article is distributed under the Creative Commons Attribution 4.0 International licence, as recorded in the arXiv licence metadata for this exact version and shown on the abstract page https://arxiv.org/abs/2601.21575v2. A full rights notice is delivered at licenses/arxiv-2601.21575-CC-BY-4.0.txt. Characters: every retained excerpt is pure ASCII, so the delivered engine, XeLaTeX with its default Latin Modern text font, reports no missing character.
	\begin{theorem}
		\label{ham_th1}
		Let $k$ be any positive integer. Then  $\alpha(Q_{8} \times C_{2}^{k}) = 16 + 2k$.
	\end{theorem}

	\begin{proof}
		Let $H=Q_8\times C_2^k$. First we observe that $\alpha(H)\geq\mu(H)=8+2k$. Again, since $\alpha(Q_8)=16$ and $\alpha(C_2^k)=2k$, it follows that $8 + 2k \leq \alpha(H) \leq 16 + 2k$. We will show that there does not exist any graph with vertex less than $16+2k$ with hamiltonian $2$-group symmetry. It is sufficient to show that $\alpha(H)\neq 8+2k, 10+2k, 12+2k, 14+2k$.\\
		First, if possible, let $\Gamma_{1}$ be a graph with $8+2k$ vertices with the automorphism group $H$. Without loss of generality we take
		\begin{align*}
			\Aut \Gamma_{1} &= \langle \sigma, \tau, \delta_{1}, \delta_{2}, \ldots, \delta_{k} \rangle \\
			&= \langle (u_{1}, u_{2}, u_{3}, u_{4})(v_{1}, v_{2}, v_{3}, v_{4}), (u_{1}, v_{3}, u_{3}, v_{1})(u_{2}, v_{2}, u_{4}, v_{4}), (a_{11}, a_{12}), \ldots, (a_{k1}, a_{k2})\rangle .
		\end{align*}
		
		Then this group has $2(2^{k} -1)+1$ elements involutions. Take $S$ be the set of all involutions. We claim that $\gamma = (u_{1}, u_{3})(u_{2}, u_{4})(a_{11}, a_{12})(a_{21}, a_{22})\ldots(a_{k1}, a_{k2})$ is also an element of $\Aut \Gamma_{1} \setminus S$, which is absurd.\\
		For $\Gamma_1$, we have the following possible edge orbits:
		\begin{align*}
			&\{[u_{i}, u_{j}]\},\ 
			\{[u_{i}, v_{j}]\},\ 
			\{[v_{i}, v_{j}]\},\\
			&\{[u_{i}, a_{l1}]\mid \forall \ l =1,2,\ldots,k\}, \
			\{[u_{i}, a_{l2}]\mid \forall \ l =1,2,\ldots,k\},\\ 
			& \{[v_{i}, a_{l1}]\mid \forall \ l =1,2,\ldots,k\},\	
			\{[v_{i}, a_{l2}]\mid \forall \ l =1,2,\ldots,k\},\\
			& \{[a_{l1}, a_{m2}]\mid \forall \ l,m =1,2,\ldots,k\}.
		\end{align*}	 
		Now, observe that
		\begin{align*}
			&\gamma([u_{i}, u_{j}]) = \sigma^{2}([u_{i}, u_{j}]),\\
			&\gamma([u_{i}, v_{j}]) = \sigma^{i+j}\tau ([u_{i}, v_{j}]),\\
			&\gamma([u_{i}, a_{l1}]) = \sigma^{2} \delta_{1} \delta_{2} \ldots \delta_{k} ([u_{i}, a_{l1}]) = \gamma([u_{i}, a_{l2}]),\\
			&\gamma([v_{i}, v_{j}]) = e,\\
			&\gamma([v_{i}, a_{l1}]) = \delta_{1} \delta_{2} \ldots \delta_{k} ([v_{i}, a_{l1}]) = \gamma([v_{i}, a_{l2}]),\\
			&\gamma([a_{l1}, a_{m2}]) = \delta_{1} \delta_{2} \ldots \delta_{k} ([a_{l1}, a_{m2}]).
		\end{align*}
		From the above calculations it follows that, for each edge $[u,v]\in E(\Gamma_1)$, an automorphism $\beta$ of $\Gamma_1$ exists such that $\gamma([u,v]) = \beta ([u,v])$. Which implies $\gamma\in \Aut \Gamma_1$. Therefore, $\alpha(H)\neq 8+2k$.  
		
		
		
		
		%%%%%%%%%%%%%
		%%%%%%%%%%%%
		
		
		%Next suppose that there exist a graph $\Gamma_{2}$ with  $ 12 + 2k$ number of vertices such that $\Aut \Gamma_{2} \cong Q_{8} \times C_{2}^{k}$. From Lemma $5$, $\sigma$ can't have three cycles. \\
		Next we suppose there exist a graph $\Gamma_{2}$ with $\Aut \Gamma_{2} = Q_{8} \times C_{2}^{k}$ with $8+2k+2$ number of vertices. Two cases arise:
		\begin{itemize}
			\item [(a)] If $\sigma = (u_{1}, u_{2}, u_{3}, u_{4})(v_{1}, v_{2}, v_{3}, v_{4})(b_{1}, b_{2})$ and $\delta_{l} = (a_{l1}, a_{l2})$ for each $l \in \{1,2,\ldots,k\}$.
			We claim  that  $\gamma = (u_{1}, u_{3}) (u_{2}, u_{4})(a_{11}, a_{12})\ldots(a_{k1}, a_{k2}) \in \Aut \Gamma_{2} \setminus S$, which is a contradiction.\\
			We have following possible edge orbits:
			\begin{align*}
				& \{[u_{i}, u_{j}]\},\
				\{[u_{i}, v_{j}]\},\
				\{[v_{i}, v_{j}]\},\\
				&\{[u_{i}, b_{1}]\},\
				\{[v_{i}, b_{1}]\},\
				\{[b_{1}, b_{2}]\},\\
				&\{[u_{i}, a_{l1}], \forall \quad l =1,2,\ldots,k\},\
				\{[u_{i}, a_{l2}], \forall \quad l =1,2,\ldots,k\},\\	
				&\{[v_{i}, a_{l1}], \forall \quad l =1,2,\ldots,k\},\
				\{[v_{i}, a_{l2}], \forall \quad l =1,2,\ldots,k\},\\
				&\{[a_{l1}, a_{m2}], \forall \quad l,m =1,2,\ldots,k\},\
				\{[b_{1}, a_{l1}], \forall \quad l =1,2,\ldots,k\},\\
				&\{[b_{1}, a_{l2}], \forall \quad l =1,2,\ldots,k\}.		
			\end{align*}
			
			% $\{[u_{i}, u_{j}]\}$, $\{[u_{i}, v_{j}]\}$, $\{[u_{i}, a_{l1}], \forall \quad l =1,2,\ldots,k\}$, $\{[u_{i}, a_{l2}], \forall \quad l =1,2,\ldots,k\}$, $\{[v_{i}, v_{j}]\}$, $\{\{[v_{i}, a_{l1}]\}, \forall \quad l =1,2,\ldots,k\}$,\\
			%	$\{[v_{i}, a_{l2}], \forall \quad l =1,2,\ldots,k\}$, $\{[a_{l1}, a_{m2}], \forall \quad l,m =1,2,\ldots,k\}$, $\{[u_{i}, b_{1}]\}$, $\{[v_{i}, b_{1}]\}$,\\
			%	 $\{[b_{1}, a_{l1}], \forall \quad l =1,2,\ldots,k\}$, $\{[b_{1}, a_{l2}], \forall \quad l =1,2,\ldots,k\}$,  $\{[b_{1}, b_{2}]\}$.    \\
			Observe that 
			\begin{align*}
				&\gamma([u_{i}, u_{j}]) = \sigma^{2}([u_{i}, u_{j}]) = \gamma([u_{i}, b_{1}]) = \gamma([v_{i}, b_{1}]) , \\
				& \gamma([u_{i}, v_{j}]) = \sigma^{i+j}\tau ([u_{i}, v_{j}]), \\
				&\gamma([u_{i}, a_{l1}]) = \sigma^{2} \delta_{1} \delta_{2} \ldots \delta_{k} ([u_{i}, a_{l1}]) = \gamma([u_{i}, a_{l2}]),\\
				&\gamma([v_{i}, v_{j}]) = e = \gamma([b_{1}, b_{2}]),\\
				&\gamma([v_{i}, a_{l1}]) = \delta_{1} \delta_{2} \ldots \delta_{k} ([v_{i}, a_{l1}]) = \gamma([v_{i}, a_{l2}]), \\
				&\gamma([a_{l1}, a_{m2}]) = \delta_{1} \delta_{2} \ldots \delta_{k} ([a_{l1}, a_{m2}]) = \gamma([b_{1}, a_{l1}]) = \gamma([b_{1}, a_{l2}]) .\\
			\end{align*}
			Hence $\gamma \in \Aut \Gamma_{2} \setminus S$, which is a contradiction.
			
			\item [(b)]   If $\sigma = (u_{1}, u_{2}, u_{3}, u_{4})(v_{1}, v_{2}, v_{3}, v_{4})$ and $\delta_{1} = (a_{11}, a_{12})(b_{1}, b_{2})$  and for each $l \in \{2,\ldots,k\}$ $\delta_{l} = (a_{l1}, a_{l2})$. Then in a similar way as above, we observe that\\
			$\gamma = (u_{1}, u_{3}) (u_{2}, u_{4})(a_{11}, a_{12})\ldots(a_{k1}, a_{k2})(b_{1}, b_{2}) \in \Aut \Gamma_{2} \setminus S$, a contradiction.
		\end{itemize}
		
		
		In similar ways, we observe that there does not exist any graph that has automorphism group isomorphic to $Q_{8} \times C_{2}^{k}$ with an $8+2k+4$ or $8+2k+6$ number of vertices.    
	\end{proof}

\end{document}
